% !TEX root = ../main.tex

\section{Results}

To diagnose the impact of shoal bathymetry on offshore transport and mixing, we define our output variables below. 

Offshore flow can be represented by velocity in the x-direction (offshore), which we call $u$. For a positive $u$, currents are moving offshore, and the opposite for a negative $u$. In our simulation figures, offshore flow represents flow moving towards the east (right), and onshore flow represents flow moving towards the west (left). 

\begin{figure}
    \centering
    \includegraphics[width=1.0\linewidth,height=0.4\textheight,keepaspectratio]{parts/figures/Hs15_100.png}
    \caption{Surface (z = 0) u-velocity (left) and kinetic energy (right) for the reference case. X axis represents distance offshore (km), and y axis represents distance alongshore (km).} 
    \label{fig:Hs15_100}
\end{figure}

In addition to offshore flow shown in our simulation output (Fig. \ref{fig:Hs15_100}), we also calculate kinetic energy from the velocity field, which is defined as 

\begin{equation}
    KE = \frac{1}{2} (u^2 + v^2 + w^2)
\end{equation}

where $u, v, w$ are the velocity components in the x, y, and z directions, respectively. 

To quantify the transport of heat and salt, we calculate the time-averaged offshore flux. The time-averaged flux of a tracer $\phi$ (representing temperature $T$ or salinity $S$) is defined as:
\begin{equation}
    \overline{u\phi} = \frac{1}{\Delta t} \int_{t_0}^{t_0+\Delta t} u(x,y,z,t) \phi(x,y,z,t) dt
\end{equation}
where we use an averaging period of $\Delta t = 10$ days. Furthermore, we compute the depth-averaged offshore flux as:
\begin{equation}
    \langle \overline{u\phi} \rangle_z = \frac{1}{H} \int_{-H}^0 \overline{u\phi} dz
\end{equation}
where $H$ is the local water depth.

Lastly, we analyze the turbulent kinetic energy (TKE) to investigate flow instabilities and mixing near the shoal. Using Reynolds decomposition ($u = \bar{u} + u'$), the time-averaged TKE is defined as:
\begin{equation}
    \overline{TKE} = \frac{1}{2} \left[ (\overline{u^2} - \bar{u}^2) + (\overline{v^2} - \bar{v}^2) + (\overline{w^2} - \bar{w}^2) \right]
\end{equation}
where the overline denotes a time-average.

For our reference case, as background flow moving northward encounters the shoal, it is deflected offshore upstream. In the region where the shoal is present on the shelf, the diverted flow follows along the shoal, and the velocity near the shoal is higher than the background velocity along the rest of the shelf. Comparing the upstream (toward the shoal) and downstream (past the shoal) velocity, offshore velocity concentrates upstream at the edge of the shoal, while onshore velocity is spread out further in the y direction downstream of the shoal. The assymetry between these two patterns suggests that the turbulence generated on the shoal as flow moves toward it redirects flow offshore. High kinetic energy is present on the surface in regions along the linear coastal ramp, and along the edge of the shoal in the x direction. Notably, the appearance of eddies on the surface, over the top of the shoal, suggests turbulent flow over the shoal that can mix the water column vertically and horizontally. 

\subsection{Parameter Sweep}

To better understand how shoal bathymetry alters flow dynamics, we conduct a parameter sweep of bathymetry-related variables. For the scope of this study, we analyze the effect of shoal height $H_S$ and shoal length $L_S$. Constants and our changing parameters are listed in Table \ref{tab:params}. 

\begin{table}[ht!]
  \centering
  \small
  \renewcommand{\arraystretch}{1.3}
  \caption{Simulation parameters for the parameter sweep. Bracketed parameters represent different test cases, with the reference case bolded.}
  \label{tab:params}
  \begin{tabular}{lll}   
    \toprule
    \textbf{Symbol} & \textbf{Parameter} & \textbf{Value(s)} \\
    \midrule
    $L_x$ & Domain length ($x$) & 100 km \\
    $L_y$ & Domain width ($y$) & 300 km \\
    $L_z$ & Domain depth ($z$) & 50 m \\
    $f$ & Coriolis frequency ($f$-plane) & $8.4168 \times 10^{-5}$ s$^{-1}$ \\
    $H_S$ & Shoal height & [5, 10, \textbf{15}] m \\
    $L_S$ & Shoal length ($x$-direction) & [\textbf{20}, 30, 40] km \\
    $\sigma_S$ & Shoal width ($y$-direction) & 8 km \\
    $L_{sh}$ & Shelf width ($x$-direction) & 50 km \\
    $H_{sh}$ & Shelf depth & 25 m \\
    $(V_N, V_S)$ & North--South inflow velocity & (0.1, 0.1) m s$^{-1}$ \\
    \bottomrule
  \end{tabular}
\end{table}

\subsection{Offshore Flux}


TODO: 

calculate volume transport? transport at shelf break? 
Plots for this: 

last time averaged window of offshore velocity (x-direction) on xy map for varying shoal heights and varying shoal lengths. 

time averaged window of offshore transport ($uT - uT_{east}$ and $uS - uS_{east}$) ? Try using recent cluster set from a couple days ago. 

stacked line plot of depth averaged and y averaged (over clipped window) offshore velocity (u) for varying shoal heights and varying shoal lengths. 

Should get a sense of shallower shoals leads to more offshore transport? 
Longer shoal leads to larger volume transport? 
Larger barrier for mean northward flow. 

avg volume transport as a dot plot where x axis is shoal length and y axis is volume transport. 

similar idea for shoal height? maybe some covariance. 


\subsection{Eddy Kinetic Energy}

TODO:

mean kinetic energy and eddy/turbulent kinetic energy depth averaged and time averaged plots. 

maybe look at mid-depth slices for KE? 

show clear transition of MKE to EKE/TKE as side by side panels. 




% This chapter contains examples of special elements that can be used in a LaTeX document, such as figures and block quotes. The purpose of this chapter is to demonstrate how to use these elements effectively in a document. This chapter is not meant to be a comprehensive guide to using these elements, but rather a brief introduction to their usage and a visual test that demonstrates how they look in the final document.

% \section{Figures}

% Figure~\ref{fig:example} shows an example of a figure that is small enough to fit where it is placed in the text (i.e. using the float placement specifier \texttt{h}).

% \begin{figure}[htbp]
%   \centering
%   \includegraphics[width=0.25\textwidth]{parts/percy.png}
%   \caption{An example figure that is placed inline. This dog's name is Percy, and he is a very good boy.}
%   \label{fig:example}
% \end{figure}

% Figure~\ref{fig:example2} shows an example of a figure that is rather large, so its location is manually changed so that it's placed at the top of the next page (i.e. using the float placement specifier \texttt{t}). The source file for this figure is the same as the one used in the previous example, but it appears to be scaled up due to the larger width specified in the \texttt{includegraphics} command. This figure is meant to demonstrate how a larger figure can be placed in a document, and how it can be scaled up to fit the desired width. Formatting figures in this manner can be helpful when the figure contains a lot of detail that would be difficult to see if it were scaled down to fit inline with the text.

% \begin{figure}[t]
%   \centering
%   \includegraphics[width=0.6\textwidth]{parts/percy.png}
%   \caption{An example figure that is too large to fit where it is placed in the text. The source file is the same one as in the previous example, but it appears to be scaled up due to the larger width.}
%   \label{fig:example2}
% \end{figure}

% Figures do not have to be created using "traditional" image files. EPS files (which can contain vector graphics like those made using Adobe Illustrator or Inkscape) can also be used. Figure~\ref{fig:example3} shows an example of a figure that is created using an EPS file. The source file for this figure is a simple vector graphic that consists of some simple geometries and text. This figure is meant to demonstrate how EPS files can be used to create figures in a LaTeX document, and how they can contain more advanced graphics than what can be achieved with a simple image file. EPS files can be particularly useful for creating figures that need to be scaled up or down without losing quality, as they are based on mathematical equations rather than pixels. Note that SVG files may also be used in a similar way, but they depend on the SVG-specific package \texttt{svg}.

% \begin{figure}[htbp]
%   \centering
%   \includegraphics[width=0.5\textwidth]{parts/vector.eps}
%   \caption{An example figure that is created using an EPS file. The source file for this figure is a simple vector graphic that consists of some simple geometries and text. Note that the quality of this figure does not degrade when it is scaled up or down, and certain features like text can be interacted with in a way that is not possible with a simple image file.}
%   \label{fig:example3}
% \end{figure}

% In some cases, it may be desirable to have figures that are wider than they are tall - or perhaps, a figure that consists of multiple subfigures. In these cases, landscape blocks and subfigures can be used. Figure~\ref{fig:example4} shows an example of a figure in a landscape-oriented page. Each subfigure is labeled with a letter (a, b, c, etc.) and has its own caption, while the overall figure has a main caption that describes the entire figure. By using separate files for each subfigure, it is possible to create a figure that contains multiple images or graphics that are related to each other, while still maintaining a clear and organized layout - all without the use of image editing software to manually combine the images into a single file. This can be particularly useful when the subfigures are generated from different sources or need to be updated separately, as it allows for greater flexibility and ease of maintenance. Additionally, using subfigures can help to improve the readability and clarity of the overall figure, as it allows for a more structured presentation of the information being conveyed.

% \afterpage{%
% \clearpage
% \begin{landscape}
% 	\begin{figure}[htbp]
% 		\centering
% 		\begin{subfigure}{0.3\textwidth}
% 			\includegraphics[width=\textwidth]{parts/percy.png}
% 			\caption{Subfigure A}
% 		\label{fig:sub1}
% 		\end{subfigure}
% 		\hfill
% 		\begin{subfigure}{0.3\textwidth}
% 			\includegraphics[width=\textwidth]{parts/percy.png}
% 			\caption{Subfigure B}
% 			\label{fig:sub2}
% 		\end{subfigure}
% 		\hfill
% 		\begin{subfigure}{0.3\textwidth}
% 			\includegraphics[width=\textwidth]{parts/percy.png}
% 			\caption{Subfigure C}
% 			\label{fig:sub3}
% 		\end{subfigure}
% 		\caption{An example of a sideways figure with three subfigures. Each subfigure is labeled with a letter and has its own caption, while the overall figure has a main caption that describes the entire figure. The source file for each subfigure is the same; this figure is only meant to demonstrate the layout of a figure with multiple subfigures, so the content of each subfigure is not important.}
% 		\label{fig:example4}
% 	\end{figure}
% \end{landscape}
% \clearpage
% }

% Notice that, for Fig.~\ref{fig:example4}, the \texttt{afterpage} command is used to ensure that the landscape figure appears on a new page, and that the \texttt{clearpage} command is used before and after the landscape block. This is necessary to ensure that the text that follows the figure is not interrupted by the figure, and that the figure is properly placed in the document. This is crucial, as the use of the landscape block without the \texttt{afterpage} and \texttt{clearpage} commands can lead to awkward appearances (e.g. large blank spaces between text and figure).

% \section{Tables}

% Tables can be created using the \texttt{tabular} environment, which allows for the creation of simple tables with rows and columns. However, for more complex tables that require features such as multi-row or multi-column cells, the \texttt{multirow} and \texttt{multicol} packages can be used. Table~\ref{tbl:example} shows an example of a table that uses these features to create a more complex layout. The content of the table is just some placeholder text, and it is meant to demonstrate how a table with multi-row and multi-column cells would look in the final document.

% \begin{table}[htbp]
% 	\centering
% 	\caption{An example of a table that uses the multirow and multicol packages to create a more complex layout. The content of the table is just some placeholder text, and it is meant to demonstrate how a table with multi-row and multi-column cells would look in the final document.}
% 	\label{tbl:example}
% 	\begin{tabular}{|c|c|c|}
% 		\hline
% 		\multirow{2}{*}{Row 1} & \multicolumn{2}{c|}{Row 1, Columns 2 and 3} \\
% 		\cline{2-3}
% 		& Column 2 & Column 3 \\
% 		\hline
% 		Row 2, Column 1 & Row 2, Column 2 & Row 2, Column 3 \\
% 		\hline
% 	\end{tabular}
% \end{table}

% Table~\ref{tbl:exampleWide} shows an example of a table defined using the \texttt{longtable} environment, which allows for tables that span multiple pages, and a landscape-oriented page. The content of the table is just some placeholder text, and it is meant to demonstrate how a long table that spans multiple pages and is oriented in landscape mode would look in the final document. Note that the \texttt{afterpage} and \texttt{clearpage} commands are used in the same way as they are for the landscape figure shown in Fig.~\ref{fig:example4}, to ensure that the table is properly placed in the document and does not interrupt the flow of the text.

% \afterpage{%
% \clearpage
% \begin{landscape}
% 	{\singlespacing
% 	\begin{longtable}{lll}
% 		\caption{An example of a table defined using the longtable environment, which allows for tables that span multiple pages, and a landscape-oriented page. The content of the table is just some placeholder text, and it is meant to demonstrate how a long table that spans multiple pages and is oriented in landscape mode would look in the final document.}\\
% 		\label{tbl:exampleWide}\\
% 		\toprule
% 		Row 1, Column 1 & Row 1, Column 2 & Row 1, Column 3 \\
% 		\midrule
% 		Row 2, Column 1 & Row 2, Column 2 & Row 2, Column 3 \\
% 		\midrule
% 		Row 3, Column 1 & Row 3, Column 2 & Row 3, Column 3 \\
% 		\midrule
% 		Row 4, Column 1 & Row 4, Column 2 & Row 4, Column 3 \\
% 		\bottomrule
% 	\end{longtable}
% 	}
% \end{landscape}
% \clearpage
% }

% \section{Special text elements}

% The following is an example of a block quote, which is indented on both sides and has a smaller font size than the main text. The content of the block quote is just some placeholder text, and it also includes a citation to show how that would look in a block quote.

% \begin{quote}
% 	Ut wisi enim ad minim veniam, quis nostrud exerci tation ullamcorper suscipit lobortis nisl ut aliquip ex ea commodo consequat. Duis autem vel eum iriure dolor in hendrerit in vulputate velit esse molestie consequat, vel illum dolore eu feugiat nulla facilisis at vero eros et accumsan et iusto odio dignissim qui blandit praesent luptatum zzril delenit augue duis dolore te feugait nulla facilisi. \cite{Example19} 
% \end{quote}

% Instead of textual statements, mathematical or chemical equations can also be expressed in a block format. Equation~\ref{eq:example} is an example of a mathematical equation that is formatted using the \texttt{equation} environment. The content of the equation is just some placeholder text, and it is meant to demonstrate how a mathematical equation would look in the final document.

% \begin{equation}
% 	\label{eq:example}
% 	E = mc^2
% \end{equation}

% Note that the equation is numbered, and the number is automatically generated by LaTeX. This allows for easy referencing of the equation in the text, as well as consistent formatting throughout the document. Additionally, using the \texttt{equation} environment allows for more advanced features such as alignment and spacing control, which can be useful for more complex equations. For example, Eq.~\ref{eq:exampleMatrix} shows an example of a matrix equation.

% \begin{equation}
% 	\label{eq:exampleMatrix}
% 	\begin{bmatrix}
% 		a & b \\
% 		c & d
% 	\end{bmatrix}
% 	\begin{bmatrix}
% 		e \\
% 		f
% 	\end{bmatrix}
% 	=
% 	\begin{bmatrix}
% 		ae + bf \\
% 		ce + df
% 	\end{bmatrix}
% \end{equation}

% An example of a theorem is shown in Theorem~\ref{thm:exampleKey}. Theorems are typically used in mathematical or scientific documents to state a proposition that can be proven based on previously established results. However, it could also be used in other types of documents to highlight a key point or idea that is being presented. The content of the theorem is just some placeholder text, and it is meant to demonstrate how a theorem would look in the final document. 

% \newtheorem{thm:example}{Theorem}
% \begin{thm:example}
% 	\label{thm:exampleKey}
% 	This is an example of a theorem. The content of the theorem is just some placeholder text, and it is meant to demonstrate how a theorem would look in the final document.
% \end{thm:example}

% Note that this document already defines custom theorems called \texttt{hyp} and \texttt{subhyp} in the preamble. These theorems are already called "Hypothesis" and "Subhypothesis", respectively, and they are meant to be used in a way that is similar to the example theorem shown above. The main difference is that the \texttt{hyp} and \texttt{subhyp} theorems are numbered in a specific way (i.e. Hypothesis 1, Subhypothesis 1a, etc.) to indicate their relationship to each other, while the example theorem shown above is not numbered at all.

% Rather than quotes, equations, or theorems, it may also be useful to express text in the form of code or pseudocode. Algorithm~\ref{alg:cap} is an example of a code block, which is formatted using the \texttt{algorithmic} environment. The content of the code block is just some placeholder text, and it is meant to demonstrate how a code block would look in the final document. 

% \begin{algorithm}[tbp]
% 	\caption{An algorithm with caption}\label{alg:cap}
% 	\begin{algorithmic}
% 		\STATE $i\gets 10$
% 		\IF {$i\geq 5$} 
% 		\STATE $i\gets i-1$
% 		\ELSE
% 		\IF {$i\leq 3$}
% 			\STATE $i\gets i+2$
% 		\ENDIF
% 		\ENDIF 
% 	\end{algorithmic}
% \end{algorithm}
